\def\topic{languages}
\ifdefined\MyHeaderIncluded
\else
\def\qwertyflag{true}

\ifdefined\MyHeaderIncluded

\else
  \newcommand{\foobar}{FooBar}
  \def\MyHeaderIncluded{true}
% if is till the end of the file

\documentclass[10pt,a4paper]{article}
\usepackage[utf8]{inputenc}
%\usepackage[german]{babel}
%\usepackage[cyrillic]{babel}
\usepackage[T1]{fontenc}
%\usepackage[T2A]{fontenc}
%\usepackage[english, russian]{babel}
\usepackage{amsmath}
\usepackage{amsfonts}
\usepackage{amssymb}
\usepackage{amsthm}
\usepackage{mathtools} %coloneqq

\usepackage{bussproofs} %

%\renewcommand{\thesubsection}{\arabic{subsection}}
%\newtheorem{thm}{Theorem}[section]
%\renewcommand{\thethm}{\arabic{subsection}.\arabic{thm}}

\usepackage{thmtools}
\usepackage{hyperref}
\hypersetup{
    colorlinks=true,
    linkcolor=blue,
    filecolor=magenta,      
    urlcolor=cyan,
    pdftitle={Overleaf Example},
    %pdfpagemode=FullScreen,
    }
%\usepackage{stix-type1}%\turnediota
\theoremstyle{definition}

\newtheorem{constructor}{Constructor}[section]
\newtheorem{iaxiom}{Axiom}
\newtheorem{axiom}{Axiom}
\newtheorem*{axiom*}{Axiom}
\newtheorem{definition}{Definition}[section]
\newtheorem*{definition*}{Definition}
\newtheorem{remark}{Remark}[section]
\newtheorem*{remark*}{Remark}
\newtheorem{recall}{Recall}[section]
\newtheorem*{recall*}{Recall}

\newtheorem{lemma}{Lemma}[section]
\newtheorem{theorem}{Theorem}[section]
\newtheorem*{theorem*}{Theorem}
\newtheorem{note}{Note}[section]
\usepackage{etoolbox}

%\newcommand*{\@Dashv}[2]{%
%  \reflectbox{$\m@th#1#2$}%
%}
\newcommand*{\separation}[3]{% 1)binder, 2)set, 3)condition
  \{#1\in #2: #3\}
}

% https://www.freecodecamp.org/news/simple-regex-tricks-for-beginners-3acb3fa257cb/
%  \\mathcal\{T\} |---> \topology
%  \\mathcal\{([^\}]*)\}   |--->  $1
\newcommand{\AND}{\ \ \&\ \ }
\newcommand{\any}{\mathrm{\ any}}
\newcommand{\borel}{\mathcal{B}}
\newcommand{\can}{\mathrm{can}}
\newcommand{\canonize}{\mathrm{canonize}}
\newcommand{\closure}{\mbox{cl}}
\newcommand{\complex}{\mathbb{C}}
\newcommand{\converges}{\longrightarrow}
\newcommand{\countable}{\mbox{\ countable}}
\newcommand{\ldefined}{\mbox{\ defined}}
\newcommand{\dom}{\mbox{dom}}
\newcommand{\Dom}{\mbox{Dom}}
\newcommand{\Done}{\noindent \textbf{Done.}\\}
\newcommand{\lexists}{\mbox{\ exists}}
\newcommand{\lforall}{\mbox{\ forall}}
\newcommand{\lunique}{\mbox{\ unique}}
\newcommand{\family}{\mathrm{\ family}}
\newcommand{\Fm}{\mathrm{Fm}}
\newcommand{\FuncSymb}{\mathrm{FuncSymb}}
\newcommand{\function}{\mathrm{\ function}}
\newcommand{\functional}{\mbox{\ functional}}
\newcommand{\IFF}{\ \Longleftrightarrow\ }
\renewcommand{\iff}{\ \longleftrightarrow\ }
\newcommand{\IMPLIES}{\ \Longrightarrow\ }
\renewcommand{\implies}{\ \longrightarrow\ }
\newcommand{\Inf}{\mathrm{Inf}}
\newcommand{\injective}{\mbox{\ injective}}
\newcommand{\injection}{\mbox{\ injection}}
\newcommand{\interior}{\mbox{int}}
\newcommand{\intersection}{\bigcap}
\newcommand{\isempty}{\mathrm{\ empty}}
\newcommand{\LB}{\mathrm{LB}}
\newcommand{\nghbset}{\mathcal{N}}
\newcommand{\nonempty}{\mathrm{\ nonempty}}
\newcommand{\Ord}{\mathrm{Ord}}
\newcommand{\power}{\mathcal{P}}
\newcommand{\pr}{\mathrm{pr}}
\newcommand{\fst}{\mathrm{fst}}
\newcommand{\snd}{\mathrm{snd}}
\newcommand{\PredSymb}{\mathrm{PredSymb}}
\newcommand{\Proof}{\noindent \textbf{Proof.}\\}
\newcommand{\ran}{\mathrm{ran}}
\newcommand{\Ran}{\mathrm{Ran}}
\newcommand{\real}{\mathbb{R}}
\newcommand{\relation}{\mathrm{\ relation}}
\newcommand{\set}{\mathrm{\ set}}
\newcommand{\singleton}{\mathrm{\ singleton}}
\newcommand{\semantics}{\mathrm{semantics}}
\newcommand{\semialgebra}{\mathcal{S}}
\newcommand{\Sup}{\mathrm{Sup}}
\newcommand{\surjective}{\mbox{\ surjective\ on\ }}
\newcommand{\Term}{\mathrm{Term}}
\newcommand{\Type}{\mathrm{Type}}
\newcommand{\topology}{\mathcal{T}}
\newcommand{\total}{\mbox{\ total\ on\ }}
\newcommand{\transitive}{\mathrm{\ transitive}}
\newcommand{\union}{\bigcup}
\newcommand{\lTrue}{\mathrm{True}}
\newcommand{\lFalse}{\mathrm{False}}
\newcommand{\lUncertain}{\mathrm{Uncertain}}
\newcommand{\lif}{\mathrm{if}}
\newcommand{\lthen}{\mathrm{then}}
\newcommand{\lelse}{\mathrm{else}}
\newcommand{\lexception}{\mathrm{exception}}
\newcommand{\lundef}{\mathrm{undef}}
\newcommand{\univocal}{\mathrm{\ univocal}}
\newcommand{\univocalat}{\mathrm{\ univocal\ at\ }}
\newcommand{\UB}{\mathrm{UB}}
\newcommand{\Un}{\mathcal{U}}
\newcommand{\Var}{\mathrm{Var}}

\newcommand{\notvarphi}{\phi}
\renewcommand{\phi}{\varphi}

\newcommand{\comment}[1]{\quad\mbox{(#1)}}
% local commands
%\preto\section{\undef\blipo} %etoolbox
\preto\subsection{\undef\blipo} %etoolbox
% https://tex.stackexchange.com/questions/105684/limit-macro-newcommand-scope-to-section
\usepackage{graphicx}
\usepackage{authblk}

\setlength\parindent{0pt} % whole doc of \noindent
\usepackage{soul}
%\usepackage{mathgl}

\fi % guard
\ifdefined\topic
\else
\def\topic{No Titel}
\fi
\title{\topic}
\author[1]{Georgy Dunaev}
\begin{document}
\maketitle
\fi
\section{Programming language with dependent types}
%[OUTDATED: write a chapter on dependent types]
\section{Primitive programming metalanguage [OUTDATED]}
\subsection{Definitions}
All inductive definitions in this section are simultaneous.
Classes are types, expressions are sets.
\begin{notation}
Various relations(:::,::,:) between text strings are introduced in this subsection.
Quotation marks are ommited:\\
Many-letter words such as Kind, Type, VarType, ConstType, Term, Fm, etc denote "Kind","Type", "VarType", "ConstType", "Term", "Fm", etc.
One-letter length latin and greek words(except ``$\iota$'' and ``$\varepsilon$'') with or without subscription are variables which can be substituted with other words.\\
Some strings are considered to be inferred, if they are obtained via (inductive) definitions.\\
Arrows ``$\Implies$'' and ``$\CImplies$'' mean, that if string to the left of them is inferred , then the string to the right is also inferred.\\
Comma ``$,$'' and ``$;$'' have higher priority than arrows. If two strings $s_1$ and $s_2$ are both inferred, then the string ``$s_1 , s_2$'' and ``$s_1 ; s_2$'' are also inferred. (Here $s_1$ and $s_2$ are meta-meta-variables.)\\
Arrow ``$\rightsquigarrow$'' denotes the relation between strings, such that that if is any substrings of any inferred text coincided with the string on the left are replaced with the string on the right, then the obtained text is stated to be inferred too. %(such rules require certain metatheorems, which will be do)
\end{notation}
%, since formulas depend on terms and vice versa.
\begin{definition}[$\Kind$]\ \\
Inductive definition of kinds:
\begin{enumerate}
\item $(\Type:::\Kind).$
\item $(\Expr:::\Kind).$
\item $(\VarType:::\Kind).$
\item $(\ConstType:::\Kind).$
\item $(\OptionType:::\Kind).$
\item If $(K_1:::\Kind)$ and $(K_2:::\Kind)$, then $((K_1\cimplies K_2):::\Kind)$.
\end{enumerate}
\end{definition}
\begin{axiom}[Kinding values of functional types]\ 
\begin{enumerate}
\item If $(A::K_1)$ and $(F::(K_1\cimplies K_2))$ then $(F(A)::K_2)$.
\end{enumerate}
\end{axiom}
%\begin{definition}
%Inductive defi
%$\Term$,
%\end{definition}
%\begin{axiom}[Kinding types with $\ConstType$]\ \\
%\end{axiom}

\begin{definition}[$\ConstType$]\ \\
%Typed constants: %NO!
Inductive definition of some constant atomic types which are main entities for which the framework is developed:
\begin{enumerate}
\item $(\Term::\ConstType).$
\item $(\Fm::\ConstType).$
\item $(\Bool::\ConstType).$
\iffalse 
\item $(\alpha::\Type) \CImplies (\Const(\alpha)::\Const\Type).$
\fi
\end{enumerate}
\end{definition}
\begin{definition}[$\OptionType$]\ \\
Inductive definition of formulas:
\begin{enumerate}
\item $(\None::\OptionType).$
\item $(\alpha::\Type) \CImplies (\Some(\alpha)::\OptionType).$
\end{enumerate}
%\item $(\alpha::\Type) \CImplies (\conseq(\alpha)::\OptionType).$
\end{definition}
\begin{definition}[$\Type$]\ \\
%\tag{def:Type}
\label{def:Type}
Inductive definition of types: for any (strings) $\alpha$ and $\beta$,
\begin{enumerate}
\iftrue % ALLOWED! % constructors must always give NEW type
\item $(\alpha::\ConstType) \CImplies  (\alpha::\Type).$
\fi
%\item $(\Term::\Type).$
%\item $(\Fm::\Type).$
\item $(\alpha::\Type)\ ;\ (\beta::\Type) \CImplies ((\alpha\to \beta)::\Type).$
%\item $(\alpha::\Type) \CImplies (\Var(\alpha)::\Type).$
\item $(\alpha::\Type)\ ;\ (\beta::\Type) \CImplies ((\alpha \times \beta)::\Type).$

\end{enumerate}
\end{definition}
\begin{definition}[$\VarType$]\ \\
Inductive definition of a kind of typed variables:
\begin{enumerate}
\item $(\alpha::\Type) \CImplies (\Var(\alpha)::\Var\Type).$
\end{enumerate}
\end{definition}
%\begin{definition}[$\Symb$]\ \\
%Inductive definition of symbols:
%\begin{enumerate}
%\end{enumerate}
%\end{definition}
\begin{definition}[$\Expr$]\ \\
Inductive definition of expressions:
%Let $\alpha$ and $\beta$ be arbitrary types, then
\begin{enumerate}
\item $(\alpha::\Type)\CImplies (\var^\alpha :: \Expr)$.
\item $(\alpha::\Type)\ ;\ (\beta::\Type)\CImplies ((x^{\alpha}.s) :: \Expr)$.
\item $(\alpha::\Type)\ ;\ (\beta::\Type)\CImplies (\pr^{\alpha\times\beta}_1 :: \Expr)$.
\item $(\alpha::\Type)\ ;\ (\beta::\Type)\CImplies (\pr^{\alpha\times\beta}_2 :: \Expr)$.
\item $(a::\Expr)\ ;\ (b::\Expr)\CImplies ((a,b) :: \Expr)$.
\end{enumerate}
\end{definition}

\subsection{Definitions related to lambda-calculus}

\begin{axiom}[Typing of $\Expr$]\ \\
Typing axioms for the expressions.
Let $\alpha$ and $\beta$ be arbitrary types, then
\begin{enumerate}
\item $(\var^\alpha : \Var(\alpha) \to \alpha)$.
%$(a:\Var(\alpha)) \Implies (a:\alpha)$.
\iffalse % its unnecessary and wrong
\item $(\const^\alpha : \Const(\alpha) \to \alpha)$.
\fi
%$(a:\Const(\alpha)) \Implies (a:\alpha)$.
\item Lambda-abstraction: $(x:\Var(\alpha))\ ,\ (s:\beta) \Implies ((x.s):\alpha\to \beta)$.
\item Application: $(g:\alpha\to \beta)\ ,\ (t:\alpha) \Implies (g(t):\beta)$.
%\item $(x:\Var\Term)\ ,\ (s:\Symb) \Implies ((x.s):\Bound\Symb)$.
\item OPairIntro: 
$(a:\alpha)\ ,\ (b:\beta) \Implies ((a,b):\alpha\times\beta)$.
\item FirstProj: $(\pr^{\alpha\times\beta}_1:\alpha\times\beta\to\alpha)$.
\item SecondProj: $(\pr^{\alpha\times\beta}_2:\alpha\times\beta\to\beta)$.
\iffalse
\item OPairElim1:
$
((a,b):\alpha\times\beta)
\Implies
(a:\alpha)
$.
\item OPairElim2:
$
((a,b):\alpha\times\beta)
\Implies
(b:\beta)
$.
\fi
\end{enumerate}
\end{axiom}

\begin{definition}[$\conseq :: (\Type \cimplies \OptionType)$]\ \\
Taking the range type $\beta$ of a function type $\alpha\to\beta$.\\
The option type $\conseq(\gamma)$ is defined by recursion on $\gamma$:
\begin{enumerate}
\item If the type $\gamma$ is made out of constant type, then return $\None$.
\item If the type $\gamma$ is a function type $\alpha\to\beta$ for some $\alpha$ and $\beta$, then return $\Some(\beta)$.
\item If the type $\gamma$ is a pair type $\alpha\times\beta$ for some $\alpha$ and $\beta$, then return $\None$.
\end{enumerate}
%\[
%\conseq(\alpha\to\beta)\ceq \Some(\beta)
%\conseq(\_)\ceq \None
%\]
So the function that takes a consequent of function type is defined. Each value of it has the $\OptionType$ kind, therefore the following inference rule is admissible:
\[
(\alpha::\Type) \CImplies (\conseq(\alpha)::\OptionType),
\]
and it $\conseq(\alpha)$ is computable for every explicitly given type $\alpha$.
\end{definition}
\begin{definition}[$\TypeOf::(\Expr\cimplies\OptionType)$]\ \\ \label{def:TypeOf}
For each well-typed expression there exists one and only one type. The corresponding function $\TypeOf$ is recursively defined by induction on the argument:
\begin{enumerate}
\item In case it is $\var^\alpha(x)$, then \[
\TypeOf(\var^\alpha(x))
%\rightsquigarrow 
\ceq 
\alpha.
\]
\iffalse
\item In case it is ...., then  \[
\TypeOf(\const^\alpha(x))
\ceq 
\alpha.
\]
\fi


\item In case it is $(x^\alpha.s)$, then 
the value of $\TypeOf((x^\alpha.s))$ is defined by induction on
$\TypeOf(s)$:
\begin{enumerate}
\item In case it is $\None$, then return $\None$.
\item In case it is $\Some(\beta)$, then return $\Some(\alpha\to\beta)$.
\end{enumerate}
\item In case it is $(g(t))$, then
the value of $\TypeOf((x^\alpha.s))$ is defined by induction on
$\TypeOf(g)$:
\begin{enumerate}
\item In case it is $\None$, then return $\None$.
\item In case it is $\Some(\gamma)$, then return $\conseq(\gamma)$.
\end{enumerate}
%\item In case it is $(g(t))$, then
%\[
%\TypeOf((x^\alpha.s))%^\beta
%\ceq 
%(\alpha\to\TypeOf(s)).
%\]
%\[
%\TypeOf(g(t))
%\ceq 
%\conseq(\TypeOf(g)).
%\]
\end{enumerate}
So the function that returns a type of an expression is defined. Each value of it has the $\OptionType$ kind, therefore the following inference rule is admissible:
\[
(e::\Expr) \CImplies (\TypeOf(e)::\OptionType),
\]
and it $\TypeOf(e)$ is computable for every explicitly given expression $e$.
\end{definition}

\begin{definition}[Well-typed expression]
The expression $e$ is well-typed if computed $\TypeOf(e)$ is not $\None$.
\end{definition}

\begin{definition}[Substitution]\ \\
Partial function of simultaneous substitution of a free variable $x$ in expression $u$ with another expression $t$.
(The definition is not yet written, but it is implemented as functions ending with ``\_simsubst'' in the project \url{https://github.com/georgydunaev/set_theory_prover/blob/main/program/checker.ML}.)
\end{definition}


\begin{definition}[Transformation]\ \\
Transformations of expressions.
\begin{enumerate}
\item Alpha-conversion: $(x.M[x/v]) \rightsquigarrow (y.M[y/v])$.
\item Beta-reduction: $((x.M) N) \rightsquigarrow (M[N/x])$.
\item Proj1-reduction: $\pr_1^{\TypeOf((a,b))}((a,b)) \rightsquigarrow a$.
\item Proj2-reduction: $\pr_1^{\TypeOf((a,b))}((a,b)) \rightsquigarrow b$.
\end{enumerate}
\end{definition}
\subsection{Basic type theoretic definitions.}
\begin{definition}[$\Bool$]\ \\
Inductive definition of booleans:
\begin{enumerate}
\item $(\lFalse:\Bool).$
\item $(\lTrue:\Bool).$
\end{enumerate}
\end{definition}
\dots
\ifdefined\qwertyflag
\undef\qwertyflag
\end{document}
\fi
